\section{BNT span consequences}\label{sec:bnt_span_consequences}

\begin{theorem}[Block-diagonal local recursion and boundary closing]
    \label{thm:block_diagonal_ground_space_intersection}
    \notready
    \uses{def:ground_space_parent, def:chain_ground_space, def:l_blk_injective,
        def:mpv, lem:blockwise_boundary_compatibility_from_traces,
        lem:normalized_boundary_matrix_identities,
        lem:left_boundary_summands_mem_block_ground_spaces,
        lem:left_boundary_trace_decomposition_mem_block_ground_spaces,
        lem:block_restrictions_mem_block_ground_spaces,
        lem:block_restrictions_iff_block_ground_spaces,
        lem:block_restriction_intersection_subspace,
        lem:period_window_block_restriction_intersection,
        lem:bnt_direct_sum_block_separation_block_intersection,
        lem:bnt_direct_sum_block_separation_period_window_block_intersection,
        lem:unital_block_separating_product_span,
        lem:product_span_block_image_direct_sum,
        lem:bnt_direct_sum_unital_block_intersection_ge,
        lem:bnt_direct_sum_unital_block_intersection,
        lem:block_sum_local_ground_space,
        lem:block_diagonal_periodic_constraints_propagated_sum,
        lem:bnt_block_diagonal_periodic_constraints_propagated_sum,
        lem:bnt_block_diagonal_periodic_chain_inclusions,
        lem:bnt_block_diagonal_open_boundary_decomposition,
        lem:bnt_block_diagonal_open_boundary_matrices,
        lem:block_diagonal_boundary_closing_equality,
        lem:block_diagonal_boundary_decomposition_inclusion,
        lem:bnt_block_diagonal_boundary_decomposition_equality,
        thm:wordspan_propagation_unital}
    Let $A_1,\ldots,A_g$ be the blocks in a block-injective canonical form, and
    assume that their MPV families are pairwise different:
    \[
        j\ne k
        \quad\Longrightarrow\quad
        \exists n\ge1,\; V^{(n)}(A_j)\ne V^{(n)}(A_k).
    \]
    Let $L_0$ be a common injectivity length. In the multi-block range used in
    \cite[Theorem~2blocks.2]{PerezGarcia2007Matrix}, namely
    \[
        g\ge2,
        \qquad
        L_0\ge1,
        \qquad
        L\ge 3(g-1)(L_0+1)+1,
    \]
    Put
    \[
        S_m:=\bigoplus_{j=1}^gG_m(A_j).
    \]
    The open-segment recursion in
    \cite[Theorem~2blocks.2]{PerezGarcia2007Matrix} says that, for every
    $m\ge L$,
    \[
        \mathbb C^d\otimes S_m
        \cap
        S_m\otimes\mathbb C^d
        =
        S_{m+1}.
    \]
    Let $H_{S_L,N}$ denote the periodic parent Hamiltonian whose local
    ground space is $S_L$. The boundary-closing conclusion of the same theorem is
    \[
        \ker H_{S_L,N}
        =
        \spn\{V^{(N)}(A_j):j=1,\ldots,g\}
    \]
    for every $N\ge L+L_0$.
\end{theorem}
\begin{proof}
    \notready
    The source proof takes a vector $\psi$ in the left-hand side and writes its
    two boundary descriptions as
    \[
        \psi=\sum_{j=1}^g \alpha_j=\sum_{j=1}^g \beta_j,
        \qquad
        \alpha_j\in\mathbb C^d\otimes G_m(A_j),
        \qquad
        \beta_j\in G_m(A_j)\otimes\mathbb C^d .
    \]
    In the source notation, the two components have boundary matrices
    $C^j_{i_1}$ and $D^j_{i_{m+1}}$: the first family leaves the initial
    physical index open, and the second family leaves the final physical index
    open.
    \[
        \alpha_j
        =
        \sum_{i_1,\ldots,i_{m+1}}
        \tr\!\left(
            A^j_{i_{m+1}} C^j_{i_1}
            A^j_{i_2}\cdots A^j_{i_m}
        \right)
        |i_1\cdots i_{m+1}\rangle,
    \]
    \[
        \beta_j
        =
        \sum_{i_1,\ldots,i_{m+1}}
        \tr\!\left(
            D^j_{i_{m+1}}
            A^j_{i_1} A^j_{i_2}\cdots A^j_{i_m}
        \right)
        |i_1\cdots i_{m+1}\rangle.
    \]
    The direct-sum separation lemma used in
    \cite[Theorem~2blocks.2]{PerezGarcia2007Matrix}, together with injectivity
    of each block, gives the blockwise matrix identity
    \[
        A^j_{i_{m+1}} C^j_{i_1}
        =
        D^j_{i_{m+1}} A^j_{i_1}
    \]
    for every $j$, $i_1$, and $i_{m+1}$. For fixed $j$, set
    \[
        E^j:=\sum_a C^j_a A^{j\,\dagger}_a .
    \]
    In the normalization used in the source proof,
    $\sum_a A^j_a A^{j\,\dagger}_a=\Id$, so
    Lemma~\ref{lem:normalized_boundary_matrix_identities} gives
    \[
        D^j_b
        =
        \sum_a D^j_b A^j_a A^{j\,\dagger}_a
        =
        \sum_a A^j_b C^j_a A^{j\,\dagger}_a
        =
        A^j_b E^j .
    \]
    Hence
    \[
        A^j_b C^j_a
        =
        A^j_b E^j A^j_a
    \]
    for every $a,b$, and cyclicity of the trace gives
    \[
        \alpha_j
        =
        \sum_{i_1,\ldots,i_{m+1}}
        \tr\!\left(A^j_{i_1}\cdots A^j_{i_m}
                  A^j_{i_{m+1}}E^j\right)
        |i_1\cdots i_{m+1}\rangle
        =
        \Gamma_{m+1}^{A_j}(E^j)
        \in G_{m+1}(A_j).
    \]
    Since $\psi=\sum_j\alpha_j$, this gives
    \[
        \psi\in\bigoplus_j G_{m+1}(A_j).
    \]
    For one block the preceding implication is
    \[
        \mathbb C^d\otimes G_m(A_j)
        \cap
        G_m(A_j)\otimes\mathbb C^d
        \subseteq
        G_{m+1}(A_j).
    \]
    The reverse inclusion is blockwise:
    \[
        G_{m+1}(A_j)
        \subseteq
        \mathbb C^d\otimes G_m(A_j)
        \cap
        G_m(A_j)\otimes\mathbb C^d .
    \]
    Thus
    \[
        \mathbb C^d\otimes G_m(A_j)
        \cap
        G_m(A_j)\otimes\mathbb C^d
        =
        G_{m+1}(A_j).
    \]
    With
    \[
        S_m:=\bigoplus_{j=1}^g G_m(A_j),
    \]
    the preceding calculation is the open-segment recursion
    \[
        \mathbb C^d\otimes S_m
        \cap
        S_m\otimes\mathbb C^d
        =
        S_{m+1}
    \]
    for every $m\ge L$. For $B=\bigoplus_j\mu_j A_j$ with every
    $\mu_j\ne0$, Lemma~\ref{lem:block_sum_local_ground_space} gives
    \[
        G_L(B)=\bigvee_jG_L(A_j).
    \]
    In the range where the sums \(\bigvee_jG_m(A_j)\) are direct, this identifies
    $G_L(B)$ with the subspace denoted in the source by
    $S=\bigoplus_j\mathcal G_L^{A^j}$.
    The passage from the open-segment recursion to the periodic chain is the
    separate boundary-closing step in
    \cite[Theorem~2blocks.2]{PerezGarcia2007Matrix}. It is not the inclusion
    $G_N(A_j)\subseteq\mathcal G_{N,L}(A_j)$, which fails for a general
    open-boundary matrix crossing the periodic boundary. The source
    boundary-closing conclusion is
    \[
        \ker H_{S_L,N}
        =
        \spn\{V^{(N)}(A_j):j=1,\ldots,g\}.
    \]
\end{proof}

\begin{lemma}[Conditional reduction to block-diagonal boundary conditions]
    \label{lem:conditional_boundary_matrix_block_split_projection_span}
    \notready
    \uses{lem:block_diagonal_commutant_reduction_from_projection_span,
        lem:eval_word_block_diagonal_tensor, lem:mpv_mem_chain_ground_space,
        lem:isup_chain_ground_space_block_le_assembled,
        thm:chain_ground_space_eq_mpv_blk, lem:center_scalar}
    Let $B^i=\bigoplus_j\mu_j A_j^i$, and assume $1<L$, $N \ge L+1$, and
    $0<m$. Suppose that every virtual sector projection $P_j$ lies in the
    pulled-back span of the length-$m$ words $B^\omega$. If a boundary matrix
    $X$ commutes with all those words, then the commutant reduction gives
    \[
        X=\bigoplus_j X_j.
    \]
    Therefore the vector obtained by applying the ground-space map to $X$ lies in
    $\sum_j \mathcal G_{N,L}(A_j)$. Consequently, suppose every vector
    $\psi\in\mathcal G_{N,L}\!\left(\bigoplus_j\mu_j A_j\right)$ has a
    boundary matrix $X$ such that
    \[
        \psi=\Gamma_N^B(X),
        \qquad
        XB^\omega=B^\omega X
    \]
    for the boundary-crossing words $\omega$ to which the preceding commutant
    reduction applies. Then
    \[
        \mathcal G_{N,L}\!\left(\bigoplus_j\mu_j A_j\right)
        \subseteq \sum_j \mathcal G_{N,L}(A_j),
    \]
    and together with the already-proved forward inclusion this gives equality.
    Composing the reverse inclusion with blockwise injective uniqueness gives
    the conditional containment in the vector space generated by the basis of
    normal tensors.  This is a conditional form of the block-diagonal
    boundary-condition step in \cite[lines~2126--2128]{Cirac2021Matrix}.
\end{lemma}
% The ground-space theorem below uses the conservative length range of
% the theorem titled "Degeneracy of the ground space v2" in
% \cite{PerezGarcia2007Matrix}.

\begin{proof}
    \notready
    \uses{lem:block_diagonal_commutant_reduction_from_projection_span,
        lem:eval_word_block_diagonal_tensor, lem:mpv_mem_chain_ground_space,
        lem:isup_chain_ground_space_block_le_assembled,
        thm:chain_ground_space_eq_mpv_blk, lem:center_scalar}
    The commutant reduction makes the pulled-back boundary matrix block
    diagonal. Restricting the same word-commutation identities to a diagonal
    block and using the injectivity of that block shows, by the scalar-commutant
    lemma, that each diagonal block is scalar. Write $\Gamma_N^B$ for the
    length-$N$ boundary-to-state map of the tensor $B$. For a block-diagonal
    boundary matrix,
    \[
        \Gamma_N^{\oplus_j\mu_j A_j}\!\left(\bigoplus_j X_j\right)
        =
        \sum_j \mu_j^N \Gamma_N^{A_j}(X_j).
    \]
    If $X_j=\lambda_j I$, this is a finite sum of scalar multiples of the
    vectors $V^{(N)}(A_j)$. The reverse inclusion and equality statements are
    reformulations of this calculation, combined with the forward inclusion and
    the blockwise uniqueness theorem.
\end{proof}

\begin{lemma}[Block-diagonal boundary reduction from blockwise word span]
    \label{lem:conditional_block_split_from_product_word_span}
    \notready
    \uses{def:blockwise_product_word_span}
    Let $A^i=\bigoplus_j\mu_j A_j^i$ be in block-injective canonical form, and
    assume $1<L$, $N \ge L+1$, and $0<m$. If
    \[
        \spn\{(A_1^\omega,\ldots,A_g^\omega): |\omega|=m\}
        =
        \prod_j \MN{D_j}
    \]
    and
    every vector in
    $\mathcal G_{N,L}\!\left(\bigoplus_j \mu_j A_j\right)$ comes with a boundary
    matrix representation commuting with all length-$m$ direct-sum words, then
    \[
        \mathcal G_{N,L}\!\left(\bigoplus_j \mu_j A_j\right)
        \subseteq \sum_j \mathcal G_{N,L}(A_j).
    \]
    Combined with the forward inclusion, this gives the corresponding conditional
    equality. In the parent-ground-space formulation, the same hypotheses imply
    containment in $\spn\{V^{(N)}(A_j): j=1,\ldots,g\}$. The blockwise word span
    and boundary representation are the displayed hypotheses.
\end{lemma}

\begin{proof}
    \notready
    \uses{lem:block_projection_span_from_product_word_span,
        lem:conditional_boundary_matrix_block_split_projection_span}
    The blockwise word-span hypothesis gives
    \[
        P_j\in
        \spn\left\{
        \bigl(\bigoplus_r\mu_r A_r\bigr)^\omega:|\omega|=m
        \right\}.
    \]
    If $X$ commutes with all displayed word products, then
    $XP_j=P_jX$ for every $j$, hence $X=\bigoplus_jX_j$. The direct-sum
    boundary formula
    \[
        \Gamma_N^{\oplus_j\mu_j A_j}\!\left(\bigoplus_jX_j\right)
        =
        \sum_j\mu_j^N\Gamma_N^{A_j}(X_j)
    \]
    gives the reverse inclusion, and the blockwise uniqueness theorem gives the
    stated span containment.
\end{proof}

\begin{lemma}[Block-diagonal boundary reduction from block-separating equations]
    \label{lem:conditional_block_split_from_bnt_block_separation}
    \notready
    \uses{def:bnt, rem:word_tuple_span_top,
        def:parent_hamiltonian_ground_space_bnt, def:bnt_span}
    Let $A^i=\bigoplus_j\mu_j A_j^i$ be in block-injective canonical form, with
    $A_1,\ldots,A_g$ a basis of normal tensors.  Assume $1 < L$ and
    $N \ge L + 1$, and suppose
    that for each block $k$ there are coefficients $c_\omega^{(k)}$ such that
    \[
        \sum_{|\omega|=S} c_\omega^{(k)} A_j^\omega
        =
        \begin{cases}
            I, & j=k,\\
            0, & j\ne k.
        \end{cases}
    \]
    If every vector in
    $\mathcal G_{N,L}\!\left(\bigoplus_j \mu_j A_j\right)$ has a boundary matrix
    representation commuting with all length-$(1 + S)$ direct-sum words, then
    \[
        \mathcal G_{N,L}\!\left(\bigoplus_j \mu_j A_j\right)
        \subseteq \sum_j \mathcal G_{N,L}(A_j).
    \]
    Together with the forward inclusion, this gives the corresponding
    conditional equality. In the parent-ground-space formulation, the same
    hypotheses imply containment in $\spn\{V^{(N)}(A_j): j=1,\ldots,g\}$. This
    is the block-injective ground-space argument of
    \cite[lines~2120--2128]{Cirac2021Matrix}, where the re-growing step is
    restricted to block-diagonal boundary conditions.
\end{lemma}

\begin{proof}
    \notready
    \uses{lem:one_block_injective_of_injective,
        lem:product_word_span_from_bnt_block_separation,
        lem:conditional_block_split_from_product_word_span}
    Finite-length injectivity gives a block length at which each block has the
    required physical--virtual correspondence. In the current length-one
    specialization, Lemma~\ref{lem:one_block_injective_of_injective} supplies
    the corresponding one-block injectivity. The displayed separation equations imply
    \[
        \spn\{(A_1^\omega,\ldots,A_g^\omega):|\omega|=1+S\}
        =
        \prod_j\MN{D_j}.
    \]
    This is Lemma~\ref{lem:product_word_span_from_bnt_block_separation} with $L=1$.
    Therefore Lemma~\ref{lem:conditional_block_split_from_product_word_span}
    applies with $m=1+S$. Its conclusion is the displayed containment, and the
    equality follows after adding the forward inclusion.
\end{proof}

\begin{theorem}[Cyclic-chain containment in the BNT span]
    \label{thm:parent_ground_space_le_bnt_span}
    \notready
    \uses{def:parent_hamiltonian_ground_space_bnt, def:bnt_span,
        thm:chain_ground_space_eq_mpv_normal}
    If $A^i=\bigoplus_j\mu_j A_j^i$ is in block-injective canonical form, where
    $A_1,\ldots,A_g$ is a basis of normal tensors, assume $g\ge2$ and
    $\mu_j\ne0$ for every $j$. Let $L_0\ge1$ be a common injectivity length for
    the blocks, and suppose
    \[
        L_0<L,
        \qquad
        L\ge 3(g-1)(L_0+1)+1,
        \qquad
        N\ge L+L_0,
    \]
    and assume the cyclic boundary-closing inclusion
    \[
        \mathcal G_{N,L}\!\left(\bigoplus_j \mu_j A_j\right)
        \subseteq
        \bigvee_{j=1}^g\mathcal G_{N,L}(A_j).
    \]
    then
    \[
        \mathcal G_{N,L}\!\left(\bigoplus_j \mu_j A_j\right)
        \subseteq
        \spn\{V^{(N)}(A_j): j=1,\ldots,g\}.
    \]
\end{theorem}

\begin{proof}
    \notready
    The length hypotheses are the conservative range used in
    \cite[Theorem~2blocks.2]{PerezGarcia2007Matrix}. The shorter range stated in
    \cite[lines~2126--2128]{Cirac2021Matrix} requires the strengthened
    block-diagonal re-growing step. These hypotheses do not by themselves give
    the displayed cyclic boundary-closing inclusion; that inclusion is the
    remaining periodic step for windows crossing the chosen cut. Since $g\ge2$
    and $L_0\ge1$,
    \[
        L\ge 3(g-1)(L_0+1)+1 \ge 3L_0+4 > L_0,
        \qquad
        N\ge L+L_0\ge L+1,
    \]
    Because $A_1,\ldots,A_g$ is a basis of normal tensors, the block MPV
    families are non-redundant:
    \[
        j\ne k
        \quad\Longrightarrow\quad
        \exists n\ge1,\; V^{(n)}(A_j)\ne V^{(n)}(A_k).
    \]
    By the assumed cyclic boundary-closing inclusion,
    \[
        \mathcal G_{N,L}\!\left(\bigoplus_j \mu_j A_j\right)
        \subseteq
        \bigvee_j \mathcal G_{N,L}(A_j).
    \]
    Theorem~\ref{thm:chain_ground_space_eq_mpv_normal} gives
    \[
        \mathcal G_{N,L}(A_j)
        =
        \spn\{V^{(N)}(A_j)\}
    \]
    for every $j$. Substituting these one-block identities into the join gives
    the displayed containment.
\end{proof}

\begin{theorem}[Cyclic-chain equality with the BNT span]
    \label{thm:gs_eq_bnt_span}
    \notready
    \uses{thm:parent_ground_space_le_bnt_span, lem:bnt_mpv_mem_ground_space,
        def:parent_hamiltonian_ground_space_bnt, def:bnt_span}
    If $A^i=\bigoplus_j\mu_j A_j^i$ is in block-injective canonical form, where
    $A_1,\ldots,A_g$ is a basis of normal tensors, assume $g\ge2$ and
    $\mu_j\ne0$ for every $j$. Let $L_0\ge1$ be a common injectivity length for
    the blocks, and suppose
    \[
        L_0<L,
        \qquad
        L\ge 3(g-1)(L_0+1)+1,
        \qquad
        N\ge L+L_0,
    \]
    and assume the cyclic boundary-closing inclusion
    \[
        \mathcal G_{N,L}\!\left(\bigoplus_j \mu_j A_j\right)
        \subseteq
        \bigvee_{j=1}^g\mathcal G_{N,L}(A_j).
    \]
    then
    \[
        \mathcal G_{N,L}\!\left(\bigoplus_j \mu_j A_j\right)
        =
        \spn\{V^{(N)}(A_j): j=1,\ldots,g\}.
    \]
\end{theorem}

\begin{proof}
    \notready
    We prove the two inclusions separately. The inclusion
    \[
        \mathcal G_{N,L}\!\left(\bigoplus_j \mu_j A_j\right)
        \subseteq
        \spn\{V^{(N)}(A_j): j=1,\ldots,g\}
    \]
    is exactly Theorem~\ref{thm:parent_ground_space_le_bnt_span}. For the reverse inclusion,
    the nonzero-weight hypothesis and Lemma~\ref{lem:bnt_mpv_mem_ground_space}
    show that each vector $V^{(N)}(A_j)$ lies in
    $\mathcal G_{N,L}(\bigoplus_j \mu_j A_j)$. Since the chain ground space is a
    linear subspace, it contains every linear combination
    of such vectors, and hence contains
    $\spn\{V^{(N)}(A_j): j=1,\ldots,g\}$. Thus
    \[
        \spn\{V^{(N)}(A_j): j=1,\ldots,g\}
        \subseteq
        \mathcal G_{N,L}\!\left(\bigoplus_j \mu_j A_j\right),
    \]
    and the two spaces are equal.
\end{proof}

\noindent\emph{Remark.} A separate kernel-identification statement for these direct-sum
block families would convert the chain-ground-space statement into the
corresponding statement for $\ker(H_N)$.
